% Per-column shading for the results tables. \cellc{s}{v} prints v on a background whose
% intensity is the INTEGER s in 0..100, 100 being the best value in that column.
%
% Two design constraints. The scale must stay legible in print and in greyscale, so it runs from
% white to a mid-tone rather than to a saturated colour, and the text stays black throughout --
% white-on-dark inverts in greyscale and is the usual way these tables become unreadable. And it
% must not imply a magnitude it does not have: the intensity encodes RANK within the column, not
% distance, because the columns hold perplexities spanning four orders of magnitude alongside
% accuracies spanning twenty points.
%
% s is an integer because \numexpr cannot divide a decimal: \cellc{0.86}{...} is a compile error.
%
% Input from the preamble AFTER xcolor, which main.tex already loads with the `table' option;
% this file deliberately does not load it itself, so that there is one place where xcolor's
% options are decided and no chance of an option clash.
\definecolor{CoCurveShade}{HTML}{2F6DB5}
\newcommand{\cellc}[2]{\colorbox{CoCurveShade!\the\numexpr#1/2\relax}{\strut #2}}
